\clearpage
\section{Hayward velocity comparison before final diagnostics}
Figure~\ref{fig:haywardvelocitycheckpoint} displays the saved standard and
catalog-DD velocity models for the authoritative subset of 1,092 events and
37 physical stations. Both use 14,262 absolute picks; the DD inversion also
uses 44,718 catalog differences. The calculation contains no waveform-derived
differences. These 21 section panels are partial counterparts to Figure 10,
with the cohort, pairing, initial model, and numerical differences retained.

The standard inversion accepted 12 updates and then rejected all line-search
steps at iteration 12. Its absolute and differential residual RMS values are
0.07869 and 0.07162 s. At that iteration, 572 of 14,262 ray calculations failed
the local stopping rule. The DD inversion reached its 18 configured updates;
its final accepted maximum source-coordinate change was 2.7614 km. Thus neither
termination establishes nonlinear convergence. The histories describe
pre-update rays; independent final-model ray diagnostics remain pending.
The plotted results retain the original straight-seed ray algorithm. The
validated initializer and constrained-direction corrections have not yet
produced a completed replacement inversion.

The color limits remain fixed across sections and methods. Standard velocity
falls below the 3.5 km/s display limit at three nodes. Across the 83,895 sampled
section points, 1,136 DD-minus-standard values are below $-0.8$ km/s and
1,265 exceed $0.8$ km/s; the full difference range is
$[-1.16214,1.37544]$ km/s. Triangular colorbar extensions expose saturation.
These differences compare two regularized estimates, not error relative to
independent velocity truth.

Regenerate these panels from the saved models with
\begin{lstlisting}
python3 plot_hayward.py --velocity-only
python3 audit.py
python3 analyze.py
\end{lstlisting}
The source and output hashes, event identities, stopping histories, and
section ranges are linked by
\path{data/hayward/figure10_velocity_manifest.json}.
The separate full-figure completion record is not written by this command.

\begin{figure}[p]
\centering
\includegraphics[width=0.98\linewidth,height=0.83\textheight,keepaspectratio]{hayward/figure10_velocity.pdf}
\caption{Partial Figure 10 counterpart from the saved original field runs:
catalog-DD P velocity, absolute-data P velocity, and DD minus absolute-data
velocity at seven written section locations. Hypocenters are within 5 km
of each section. The full computational depth interval includes the layer
above the geoid needed to contain elevated receivers. Colorbar extensions
mark values outside the common display limits. These numerical outputs
retain unresolved ray and inverse-problem convergence limitations.}
\label{fig:haywardvelocitycheckpoint}
\end{figure}
